% ---------------------------------------------------------------------------
% LC-equivalent inputs, non-LC-equivalent outputs under graph fusion.
% Node T holds a triangle T1,T2,T3; nodes X in {A,B,C} each hold X1,X2,X3.
%   (a) X-parts are paths X1-X2-X3
%   (b) X-parts are triangles: tau_{X2} toggles (X1,X3).
%       So (b) = tau_{A2} tau_{B2} tau_{C2} (a).
% Fusions F(A1,T1), F(B1,T2), F(C1,T3) (green qubits, dashed pairing)
% toggle N(x) x N(t) and delete x,t:
%   (a) -> triangle A2,B2,C2 with pendants A3,B3,C3
%   (b) -> complete graph K6, LC-equivalent to GHZ_6
% The two outputs are not LC equivalent (checked by LC-orbit enumeration).
% Output vertices are drawn where they sit in the input.
% B- and C-parts point sideways, level with T2 and T3, to save height.
% gold = edge added by tau_{X2}.
% ---------------------------------------------------------------------------
{\footnotesize
% the X-part of node #1, rotated by #2 degrees around T and shifted
% vertically by #4; #3 = style of (X1,X3)
\newcommand{\lcfx}[4]{%
  \begin{scope}[shift={(0,#4)},rotate=#2]
    \node[gv,fuv] (#11) at (0,1.3) {$#11$};
    \node[gv] (#12) at (-0.42,1.95) {$#12$};
    \node[gv] (#13) at (0.42,1.95) {$#13$};
  \end{scope}
  \draw[ge](#11)--(#12);\draw[ge](#12)--(#13);\draw[#3](#11)--(#13);}
% input panel: #1 = style of the (X1,X3) edges (draw=none: absent)
\newcommand{\lcfin}[1]{%
  \begin{tikzpicture}[baseline={(0,0.5)},x=1cm,y=1cm,
      gv/.append style={minimum size=5.5mm,font=\scriptsize},
      fpair/.style={Ycol,line width=1.1pt,dash pattern=on 1.5pt off 1pt}]
    \node[gv,fuv] (T1) at (90:0.42)  {$T1$};
    \node[gv,fuv] (T2) at (210:0.42) {$T2$};
    \node[gv,fuv] (T3) at (330:0.42) {$T3$};
    \draw[ge](T1)--(T2)--(T3)--(T1);
    \lcfx{A}{0}{#1}{0}
    \lcfx{B}{90}{#1}{-0.21}
    \lcfx{C}{-90}{#1}{-0.21}
    \draw[fpair](A1)--(T1);\draw[fpair](B1)--(T2);\draw[fpair](C1)--(T3);
  \end{tikzpicture}}
% output panel: #1 = edges
\newcommand{\lcfout}[1]{%
  \begin{tikzpicture}[baseline={(0,0.5)},x=1cm,y=1cm,
      gv/.append style={minimum size=5.5mm,font=\scriptsize}]
    \foreach \X/\r/\s in {A/0/0,B/90/-0.21,C/-90/-0.21}{
      \begin{scope}[shift={(0,\s)},rotate=\r]
        \node[gv] (\X2) at (-0.42,1.95) {$\X2$};
        \node[gv] (\X3) at (0.42,1.95) {$\X3$};
      \end{scope}}
    \begin{pgfonlayer}{background}#1\end{pgfonlayer}
  \end{tikzpicture}}
% panel label, raised above the baseline
\newcommand{\lcflab}[1]{\raisebox{4mm}{\normalsize(#1)}}
\newcommand{\lcfarr}{{\small$\xrightarrow{\ \substack{\textstyle\fuse{A1,T1}\\[2pt]
    \textstyle\fuse{B1,T2}\\[2pt]\textstyle\fuse{C1,T3}}\ }$}}
\setlength{\tabcolsep}{1pt}%
\resizebox{0.85\linewidth}{!}{%
\begin{tabular}{@{}c@{}c@{}c@{}c@{}c@{}}
\lcflab{a} & \lcfin{draw=none} & \lcfarr
  & \lcfout{\draw[ge](A2)--(A3);\draw[ge](B2)--(B3);\draw[ge](C2)--(C3);
             \draw[ge](A2)--(B2)--(C2)--(A2);} & \lcflab{a'} \\\noalign{\vskip 12pt}
  & {\large$\Big\updownarrow\ \localComp{A2}\localComp{B2}\localComp{C2}$}
  & & {\Large$\not\sim_{\mathrm{LC}}$} \\\noalign{\vskip 12pt}
\lcflab{b} & \lcfin{newe} & \lcfarr
  & \lcfout{\foreach \u/\v in {A2/A3,A2/B2,A2/B3,A2/C2,A2/C3,A3/B2,A3/B3,
               A3/C2,A3/C3,B2/B3,B2/C2,B2/C3,B3/C2,B3/C3,C2/C3}
               {\draw[ge](\u)--(\v);}} & \lcflab{b'}
\end{tabular}}%
}
